\documentclass[11pt]{book}

\usepackage{graphicx}
\usepackage{eso-pic}
\usepackage{tcolorbox}
\usepackage{polyglossia}
\setdefaultlanguage[indentfirst=false,forceheadingpunctuation=false]{russian}
\setotherlanguages{english}

\RequirePackage[a4paper, headsep = 0.5 \headsep, left=2.5cm, right=2.1cm, top=2cm, bottom=2.1cm]{geometry}

\setmainfont{Times New Roman}
\newfontfamily\cyrillicfont{Times New Roman}[Script=Cyrillic, Ligatures=TeX]

\usepackage{enumitem}
\setlist[itemize, 1]{label =\raisebox{-0.3\height}{\scalebox{1.8}{\textbullet}}}

\title{Феноменология: Основы}
\author{Дэн Захави}
\date{2019}
\begin{document}

\AddToShipoutPictureBG*{\includegraphics[width=\paperwidth]{title.png}}

\null
    \thispagestyle{empty}%
    \addtocounter{page}{-1}%
    \newpage

\tableofcontents

\chapter{Введение}

Феноменология считается одной из доминирующих традиций в философии XX века. Ее основал Эдмунд Гуссерль (1859-1938)\footnote{Эдмунд Гуссерль пришел в философию в относительно позднем возрасте. Он родился в Простеёве, Моравия (тогда входившая в состав Австрийской Империи), изучал физику, математику, астрономию и философию в Лейпциге, Берлине и Вене. В 1883г он защитил докторскую диссертацию по математике, и в последующие годы серьезно заинтересовался философией , отчасти благодаря посещению лекций психолога и философа Франца Брентано. Первая крупная работа Гуссерля --- ''Логические исследования'' --- была опубликована в 1900-1901гг, и именно на ее основе он получил приглашение в Геттинген, где преподавал с 1901 по 1916 год. Следующая крупная работа --- ''Идеи к чистой феноменологии и феноменологической философии I'' --- была опубликована в 1913г. В 1916г Гуссерль переезжает во Фрайбург, где получает кафедру философии, прежде занимаемую Генрихом Риккертом. В последующие годы его ассистентами работают Эдит Штайн и Мартин Хайдеггер, который, когда Гуссерль выйдет на пенсию в 1928г, займет его место. Вскоре после выхода на пенсию Гуссерль публикует ''Формальную и трансцендентальную логику'' (1929) и ''Картезианские размышления'' (1931). В тридцатые годы Гуссерль страдал от условий, навязанных немецким национал-социалистическим режимом: лишенный какой-либо официальной академической деятельности из-за своего еврейского происхождения, он теряет право преподавать, публиковаться, и, в конечном итоге, свое немецкое гражданство. Он, однако, продолжает свою работу, в такие темные времена еще более настаивая на актуальности философии. В 1935г Гуссерля приглашают читать лекции в Вене и Праге, эти лекции в последствии лягут в основу его последней работы ''Кризис европейских наук и трансцендентальная феноменология'' (1936). Вскоре после его смерти 27 апреля 1938г молодому монаху-францисканцу Герману Лео ван Бреде удается контрабандой вывезти многочисленные исследовательские рукописи Гуссерля из Германии в монастырь в Бельгии. Перед началом Второй Мировой Войны в Левене был создан Архив Гуссерля, где до сих пор можно найти оригинальные рукописи.}, другие влиятельные сторонники --- Мартин Хайдеггер (1889-1976)\footnote{Хайдеггер начал изучать католическую теологию и средневековую философию во Фрайбурге, но в 1911г решил сосредоточиться на философии. В 1913г он защищает диссертацию, а два года спустя --- хабилитацию ''Учение Дунса Скотта о категориях и значении'', которая затем будет передана Риккерту --- философу, преемником которого был Гуссерль. Хайдеггер работает ассистентом Гуссерля в 1918-1923гг, и затем становится профессором университета в Марбурге. В 1927 году он публикует главную свою работу --- ''Бытие и время'' --- а в 1928 занимает кафедру Гуссерля во Фрайбурге. В 1929г Хайдеггер прочел свою знаменитую вступительную лекцию ''Что такое метафизика?''. После прихода к власти Гитлера Хайдеггер был избран ректором Фрайбургского университета и стал членом нацистской партии. Однако менее чем через год он уходит с этого поста и постепенно отходит от университетской политики. Как показывает публикация в 2014г его так называемых ''Черных тетрадей'', флирт Хайдеггера с нацизмом не был ни кратковременным, ни поверхностным. До 1944г он регулярно читал лекции, но после окончания войны ему запретили преподавать из-за политических убеждений, а в 1946г он был лишен звания профессора. В 1949г его восстанавливают в звании почетного профессора, и с тех пор и до своей смерти о читает много лекций и публикует такие центральные работы как ''Язык'' (1950), ''Размышления о строительстве и обитании'' (1951) и ''Вопрос технологии'' (1953).}, Жан-Поль Сартр (1905-1980)\footnote{Сартр изучал философию в Высшей нормальной школе, где встретил целое поколение ведущих французских интеллектуалов, в том числе Симону де Бовуар, Раймона Арона, Мориса Мерло-Понти, Симону Вейль, Эммануэля Мунье, Жана Ипполита и Клода Леви-Стросса. Его отношения с коллегой-философом Бовуар станут легендарными. В период 1931-1945гг Сартр преподает в средних школах Гавра, Лана и Парижа. С философией Гуссерля и Хайдеггера его знакомят в начале 1930гг Арон и Левинас, а в 1933-1934гг он отправляется в Берлин изучать феноменологию. Во второй половине тридцатых годов Сартр публикует четыре книги о сознании: ''Трансцендентность эго'', в которой рассматривается структура самосознания (1936), ''Воображение'' (1936) и ''Воображаемое'' (1940), а также ''Очерк теории эмоций'' (1939). Все четыре эти работы свидетельствуют о его знакомстве с философией Гуссерля. С началом Второй Мировой Войны Сартр был призван в армию и в 1940г попал в плен к немцам. После освобождения в 1941г он становится активным участником сопротивления и в 1943г публикует ''Бытие и ничто'', ясно демонстрирующий влияние Хайдеггера. В 1945г Сартр основывает литературно-политическое обозрение Les Temps Moderne (''Новые времена''), которое редактирует в течение многих лет (периодически совместно с Мерло-Понти). После войны Сартр решает прекратить преподавание и посвящает себя литературной и редакционной работе. Сартр, таким образом, один из очень немногих влиятельных философов XX века, который не работал в университете. В послевоенный период политическая активность Сартра возрастает, возрастают и его симпатии к марксизму и восхищение СССР, хотя членом коммунистической партии он так и не станет. Поддержка им СССР остается неизменной до вторжения в Венгрию в 1956г. В 1960г Сартр публикует свое второе основное произведение ''Критика диалектического разума'', что свидетельствует о его политической и социальной активности. В 1964г Сартру присудят Нобелевскую премию по литературе, но он из принципа откажется от нее. Он останется политически активным до позднего возраста, выступая против войны Франции в Алжире, активно выступая против войны во Вьетнаме вместе с Бертраном Расселом, а также поддерживая протесты в Париже в мае 1968г. На его похороны после смерти в апреле 1980г придет около 50тыс. человек.}, Морис Мерло-Понти (1908-1961)\footnote{Мерло-Понти изучал философию в престижной Высшей нормальной школе в Париже и был первым иностранным посетителем недавно созданного Архива Гуссерля в Левене в апреле 1939г. Свою первую книгу ''Структура поведения'' он публикует в 1942г, а свою, пожалуй, главную работу --- ''Феноменология восприятия'' --- в 1945г. В 1949г он становится профессором детской психологии и педагогики в Сорбонне, а в 1952г будет избран на кафедру философии в Коллеж де Франс, став самым молодым из когда-либо назначенных на эту должность, и будет занимать ее до своей смерти в мае 1961г. После Второй Мировой Войны Мерло-Понти все активнее занимается политикой и публикует ряд книг с эссе, среди которых ''Гуманизм и террор'' (1947), ''Смысл и бессмыслица'' (1948) и ''Приключения диалектики'' (1955). Параллельно с этими политическими интересами он продолжает преподавать, и многие из его лекций в Сорбонне и Коллеж де Франс свидетельствуют об обширном интересе к другим дисциплинам, включая детскую психологию, структурную лингвистику, этнологию и психоанализ. В 1960г выходит еще один том ''Знаков'', состоящий из эссе, а в 1964г будет посмертно опубликован отрывок ''Видимое и невидимое'', которое многие считают вторым главным произведением Мерло-Понти.} и Эммануэль Левинас (1906-1995)\footnote{Эммануэль Левинас родился в Каунасе в Литве (тогда входившей в состав Российской Империи). После окончания средней школы он уезжает во Францию, где начинает изучать философию в Страсбурге. В 1928г он обучается у Гуссерля и Хайдеггера во Фрайбурге, и затем совместно с Г.Пайффером переведет ''Картезианские размышления'' Гуссерля на французский. С публикацией в 1931г его докторской диссертации ''Теория интуиции в феноменологии Гуссерля'' за ним устанавливается репутация как одного из ведущих специалистов по немецкой феноменологии во Франции. К моменту начала Второй Мировой Войны Левинас стал гражданином Франции и был призван в армию. Как и Сартр, он вскоре попал в плен к немцам, где провел остаток войны, благодаря чему избежал участи большей части своей еврейской семьи, убитой немцами. После войны Левинас становится директором École Normale Israélite Orientale, а затем получает должности в университетах Пуатье (1961), Нантера (1967), а с 1973г --- в Сорбонне в Париже. Ранние послевоенные работы Левинаса ''От существования к существующему'' (1947), ''Время и другой'' (1948) и ''Открытие существования с Гуссерлем и Хайдеггером'' (1949) все еще демонстрируют его долг как перед Гуссерлем, так и перед Хайдеггером, но уже рассматривают характерные для собственного мышления Левинаса темы, такие как роль другого и отношение между этикой и онтологией. Рассмотрение этих тем достигает своей кульминации в ''Тотальность и бесконечность'' (1961) и ''Иное, чем бытие или за пределами сущности'' (1974). Помимо множества философских книг, Левинас опубликует также серию талмудических комментариев и эссе об иудаизме (в том числе ''Трудная свобода'' (1963) и ''Девять талмудических чтений'' (1968 / 1977)).}. Влияние феноменологии обусловлено в том числе возможностью понимания почти все последующих теоретических формирований в немецкой и французской философии либо как расширение феноменологии, либо как реакцию на нее. Понимание феноменологии, значит, необходимо для понимания последующего развития философии XX века.

За прошедшие годы феноменология сделала значительный вклад во многие области философии и предложила новаторский анализ интенциональности, восприятия, воплощения, эмоций, самосознания, интерсубъективности, темпоральности, историчности и истины. Она целенаправлено критикует редукционизм, объективизм и сциентизма. Подробно рассматривая человеческое существование и понимая субъекта как воплощенное и социально и культурно встроенное бытие-в-мире, феноменология делает решающий вклад в такие эмпирические дисциплины как психиатрия, социология, психология, литературоведение, антропология и архитектура.

Многие из ныне классических работ по феноменологии были написаны в первой половине XX века, но в последние годы интерес к ней возобновился. Феноменология сейчас переживает своего рода ренессанс.

Феноменология развивалась как во многом разнородное движение с множеством ответвлений, все феноменологи после Гуссерля дистанцировались от различных аспектов его исходной программы, и вряд ли корректно говорить о феноменологии как о философской системе с четко очерченным корпусом доктрин, однако не следует упускать из виду те всеобъемлющие проблемы и темы, которые объединяли и продолжают объединять ее сторонников.

В введении, конечно, невозможно отдать должное всем феноменологам, и я буду опираться прежде всего на работы Гуссерля, раннего Хайдеггера и Мерло-Понти --- трех мыслителей, чье влияние на развитие феноменологии бесспорно. Многие вводные книги по феноменологии содержат отдельные главы, посвященные этим мыслителям, но я поступлю иначе: вместо указания на их различия --- часто, на мой взгляд, преувеличивающиеся из-за фундаментально ошибочного толкования идей Гуссерля --- я сосредоточусь на том, что в них общего.

В первой части книги рассмотрена концепция философии согласно феноменологии --- вопрос метода, фокус на виде от первого лица, анализ проживаемого мира --- также развитие феноменологической традиции.

Во второй части я откладываю в сторону методологические соображения чтобы рассмотреть несколько примеров и моделей феноменологического анализа. Я рассмотрю феноменологические исследования пространственности и воплощенности, а также анализ интерсубъективности и сообщества.

Третья часть книги демонстрирует применение феноменологии вне философии. В частности, я опишу ее влияние на социологию, психологию и когнитивистику.

Главная цель этой книги --- не защита феноменологии, но ее представление: изложение ряда ее идей в максимально доступной форме без защиты их от разного рода критики. Я старался как можно меньше внимания уделять интерпретационным разногласиям и противоречиям. Так что если вы ищете систематическую защиту феноменологии и широкое соотнесение ее положений и выводов с исследовательской литературой, эта книга вам не подойдет\footnote{См., например, Захави 2003, 2005, 2014, 2017 и Галлахер и Захави 2012.}.

\part{Фундаментальные вопросы}

Следующие шесть глав представляют базовый обзор феноменологии: какого рода исследованиями она занимается, сосредоточена ли он преимущественно или даже исключительно на разуме или же в равной степени касается и мира, что такое феномен и как можно его исследовать, какие есть способы взаимодействия с миром, в чем разница между разговором о дереве и этого дерева восприятием, как мир науки соотносится со знакомым из повседневного опыта миром, что значит причастность феноменологии к трансцендентальной философии. В главах будут представлены некоторые глобальные темы и проблемы феноменологии, описаны ее методы и очерчено ее развитие.

\chapter{Феномен}

Строго говоря, феноменология занимается исследованием феноменов. Но что такое феномен, и какие феномены интересны феноменологам? Возможно, их занимают удивительные, впечатляющие, феноменальные явления? В своей автобиографии Симона де Бовуар рассказывает о первом знакомстве Сартра с феноменологией. Они были в коктейль-баре с Раймоном Ароном, только что вернувшимся из Германии, и Арон указал на заказанный им абрикосовый коктейль и сказал: ''Видите ли, мой дорогой друг, если вы феноменолог, то можете заниматься философией, говоря об этом коктейле!''. Арон был совершенно прав: даже повседневный опыт простых объектов может служить отправной точкой для феноменологического анализа. Если философия хочет избежать тупика устаревших абстракций, ей необходимо воссоединиться с богатством повседневной жизни. Впрочем, стоит отметить нацеленность феноменологии на скорее вопросы ''как'' чем ''что'' касательно объектов. Вместо того, чтобы говорить о, скажем, весе, редкости или химическом составе объекта, феноменология занимается его проявлениями, тем как он нам является. Есть существенные отличия в том, как представляют себя посуда, произведение искусства, мелодия, положение дел, число и человек. Более того, один и тот же объект может проявлять себя по-разному: с той или иной точки зрения, при сильном или слабом освещении, как воспринимаемый, воображаемый, желаемый, избегаемый, предвкушаемый или вспоминаемый. Феноменологию можно рассматривать как философский анализ различных типов данности.

\section{Будильник}

Все это звучит, наверное, абстрактно и даже немного эзотерично, так что давайте рассмотрим конкретный пример. В винтажном магазине в центре Копенгагена я ищу подарок на день рождения своему другу, и замечаю старинный будильник. Как будильник проявляет себя, что это за явление? Для начала следует признать, что простого ответа нет, поскольку будильник может проявлять себя по-разному. Я могу его не только увидеть, потрогать и услышать его, но также представить, как если бы он был дан мне на картинке, или просто-таки использовать. Для простоты давайте сосредоточимся на том, как будильник дан в рассматриваемой ситуации, т.е. как он дан в восприятии. То, как будильник будет выглядеть, зависит от освещенности --- будет это естественный солнечный свет, неоновые огни, прожекторы или что-то еще. Впрочем, вне зависимости от обстоятельств и даже при оптимальном освещении я никогда не смогу увидеть его целиком, он всегда будет дан с некой перспективы: если я посмотрю на будильник когда он стоит на столе в магазине, то смогу увидеть его верхнюю и боковые части, но не смогу увидеть заднюю, нижнюю или внутреннюю. Если же я пойду вокруг стола, то смогу увидеть обратную сторону будильника, и если подниму его со стола, то смогу осмотреть и нижнюю, но он всегда будет дан с перспективы: раскрытие одних аспектов будет сопровождаться исчезновением с поля зрения других. Это все может показаться довольно тривиальным наблюдением, повседневным и не требующим никаких дальнейших размышлений, но в руках феноменологов оно содержит ключ к множеству проблем. Прежде всего следует отметить нашу уверенность в наличие у будильника нечта большего чем то что нам о нем дано, хотя мы никогда не видим его целиком: есть его части и свойства, отсутствующие в нашем восприятии. Эти аспекты --- даже не будучи даны в восприятии --- играют в нем важную роль: воспринимая объект, мы переживаем больше, чем интуитивно данное, так что видимая сторона будильника указывает на другие, с этой перспективы не данные но могущие себя обнаружить при дальнейшем исследовании. Видимое, значит, никогда не дано изолированно, а окружено неким горизонтом, участвующем в формировании его значения.

Горизонт этот охватывает больше чем просто невидимые на мгновение аспекты рассматриваемого объекта: мы никогда не сталкиваемся с изолированными объектами, только с включенными в более широкий контекст. Так, будильник, на который я смотрю, стоит на столе в определенной комнате, в зависимости от назначения которой (торговый зал, аудитория или кабинет) будет являться по-разному, с разными значениями.

Поле моего сознания не исчерпывается будильником даже когда я наблюдаю за ним. Он может быть окружен другими часами, чашками, ручками, книгами и т.д., и сосредотачиваясь на будильнике я не буду обращать внимания на все это, но я и не буду вовсе забывать о нем, так же как и о полу, на котором стою, об освещенности комнаты и т.д.: оно будет дано как основа, контекст, части целостности фона будильника. И будучи лишь фоном они все еще могут стать объектом внимания при его смене. Возможность такого тематического изменения основана на со-данности поля основной теме так что в этом поле возможно мысленное перемещение.

Что значит данность с перспективы? В восприятии будильник дан всегда под определенным углом и на определенном расстоянии от воспринимающего. Но что это говорит нам о воспринимающем? Для такой данности будильника воспринимающий должен находиться в том же, что и будильник, пространстве, а расположенность в пространстве значит воплощенность: бестелесный воспринимающий не имел бы места в пространстве, так что будильник может быть дан только как данный воплощенному воспринимающему, и нет чисто интеллектуальной перспективы или взгляда из неоткуда --- есть только воплощенная.

Решающая роль тела в восприятии может быть проиллюстрирована и несколько иначе. Даже всегда изначально сталкиваясь с очень ограниченным взглядом на объект, мы редко остаемся таким взглядом удовлетворены. Как указывает Гуссерль, объект манит нас к дальнейшему исследованию:

\smallskip
\leftskip1cm\relax
\rightskip1cm\relax

Здесь еще есть на что посмотреть, поверни и увидь все стороны, позволь взгляду пробежаться по мне, приблизься ко мне, открой меня, раздели меня, продолжай смотреть снова и снова, поворачивай чтобы увидеть все стороны: так ты узнаешь меня, все чем я являюсь, все мои качества.

\leftskip0cm\relax
\rightskip0cm\relax
\smallskip

Как провести такое исследование, как лучше узнать будильник? Взяв его в руки и повернув, или обойдя стол так чтобы рассмотреть обратную его сторону: все это требует телесного участия и взаимодействия, предполагает их. Такое перцептивное исследование, значит --- телесная деятельность: мы двигаем глазами, головой, туловищем, руками и вообще всем телом. Эта деятельность не мгновенна, она требует времени, время играет решающую роль: если сначала посмотреть на переднюю часть будильника а затем осмотреть заднюю его часть, то передняя постепенно исчезнет из поля зрения, но все еще будет осознаваться. Такое накопление знаний о будильнике возможно именно потому что возможно сохранение виденного в прошлом. Меняя перспективу и положение, мы не воспринимаем сначала переднюю часть, затем боковую и затем заднюю, как если бы рассматривали три разных отдельных снимка. Взяв его в руки и покрутив, мы наблюдаем постепенную смену его данности, и постепенность эта значит что и наш поток сознания не серия мгновенных обособленных восприятий и имеет особую целостную временную и эмпирическую структуру. Время играет и другую роль: подчеркивая важность контекста и горизонта, о них следует думать не только в пространственных, но и во временных терминах. Мы сталкиваемся с настоящим на основе прошлого и планами и ожиданиями касательно будущего, прошлых опыт не теряется и не оставляет нас нетронутыми. Так, будильник в магазине заинтересовал меня не только потому что по опыту знаю что мой друг крепко спит, но и потому что намереваюсь сделать ему подарок.

Это подводит нас к еще одному замечанию: будильник дан мне не как приватный, не как данный лишь мне --- это публичный объект, который и другие могут наблюдать и использовать, и потому я подумал о его покупке. Даже будучи явлен мне лишь частично, он может в то же время быть явлен другим другими своими аспектами, мне недоступными.

\section{Данность и действительность}

Какое значение это все имеет? Обращая внимание на данность различных объектов, феноменология говорит лишь о явлениях, т.е. об \textit{очевидной} части природы объектов. Целью же науки обыкновенно понимается постижение объектов такими, какие они действительно есть.

В большей части философской традиции феномен определялся как представленность нам некоторого объекта, его видимость и мыслимость, и противопоставлялся объекту как он на самом деле есть. И, значит, в стремлении открыть и определить каков объект на самом деле следует выйти за рамки лишь феноменального. Понимай феноменология феномен так, она бы изучала лишь субъективное, кажущееся, \textit{поверхностное}. Как указывает Хайдеггер в седьмом разделе ''Бытия и времени'', феноменология опирается на совершенно иное, более классическое понимание феномена как того что себя проявляет, раскрывает. Феноменология, таким образом, не о лишь кажущемся, как Хайдеггер указывал также в курсе лекций, прочитанном им за несколько лет до написания ''Бытия и времени'':

\smallskip
\leftskip1cm\relax
\rightskip1cm\relax

Феноменологически абсурдно говорить о явлении как о скрывающем за собой нечто из чего оно как его представляющее и выражающее. Феномен это то за чем ничего более нет, и даже некорректно пытаться о чем-то за ним говорить потому как феноменологическая данность уже самодостаточна.

\leftskip0cm\relax
\rightskip0cm\relax
\smallskip

Некоторые могут утверждать, что феномен --- нечто субъективное, завесой дыма скрывающее объективную реальность. Феноменологи отвергают так называемую доктрину двух миров о должности принципиального различия между представленным и могущим быть понятым нами миром и миром каков он на самом деле есть. Это, конечно, не значит отрицание различия между лишь видимостью и реальностью --- в конце концов, некоторые видимости вводят в заблуждение --- но различие это не между двумя отдельными сферами, а между двумя способами явленности, между данностью при поверхностном взгляде и в результате, например, тщательного научного исследования. Феноменологи обыкновенно утверждают обладание явленным нам миром --- будь то в восприятии, повседневности или научном анализе --- всей необходимой реальностью и объективностью. Существование скрытого, закулисного и превосходящего всякую данность, мира отвергается как пустая спекуляция. Само подобное предложение, настаивают феноменологи, содержит категориальную ошибку, некорректное применение понятия реальности и злоупотребление им: вместо ее определения как недоступной и непостижимой, следует локализовать объективность внутри видимого мира.

Итак, какой же вывод следует сделать касательно области феноменологического исследования? Что поставлено на карту в феноменологическом анализе и следует ли эти исследования понимать как исследования сознания? Как теперь должно быть ясно, феноменология не фокусируется в первую очередь или даже отдельно на структуре разума: надлежащим фокусом феноменологического анализа является диада разум -- мир, или, как в конечном счете станет видно, триада я -- другой -- мир.

\chapter{Интенциональность}

Я могу думать о простых числах или далеких галактиках, представлять единорогов, чувствовать вкус ромашкового чая, бояться дементоров, вспоминать события из своего детства или узнавать датскую королеву на фотографии. Об эмпирической жизни, значит, фундаментально ошибочно думать как лишь о смеси более или менее интенсивных внутренних ощущений и чувственных состояний. Видение, слышание, воспоминание, воображение, мышление, ненависть и страх всегда о чем-то, всегда на что-то направлены --- сознание не самодавлеюще, а, скорее, самопреодолевающе. Феноменологи используют термин ''интенциональность'' для обозначения ''указания за пределы себя'', присущего сознанию.

Хотя понятие интенциональности имеет долгую историю, восходящую по крайней мере к Аристотелю, и играло центральную роль в средневековой философии, оно подвергнуто серьезному анализу в первой крупной работе по феноменологии --- ''Логические исследования'' Гуссерля (1900-1901). Гуссерль продолжал работу над этой темой до конца своей жизни, и его исследование позволило прояснить связь и различие между эмпирической субъективностью и объектами мира.

Будучи интенционально, сознание не замкнуто в себе, а занято объектами и событиями совершенно по своей природе от него отличными. Вспоминаемый мной праздничный торт, например, сильно отличается от его восприятия: он весит 3 килограмма и его можно есть, тогда как соответствующий ему опыт ничего не весит и его не получится съесть. Торт, к тому же, не о чем-то, ни на что не направлен, в отличие от восприятия.

Интенциональность аспектна и перспективна. Объект никогда не осознается как таковой, всегда определенным образом, с точки зрения или через описание. О смартфоне можно думать как о средстве связи, как о подарке, как об эффективном способе хранения музыки или как об источнике раздражения (когда он не работает должным образом). Но помимо представления может варьироваться форма представления: вместо того чтобы смотреть не таблицу, можно вообразить ее, судить о ней, запомнить ее и т.д.. Один на один и тот же объект сознание может быть направлено по-разному и одним и тем же образом оно может быть направлено на разные объекты. Например, можно сомневаться в самых разных вещах: что финансовый кризис продолжится, что выборы были честными, что климатический кризис выдуман, и можно видеть, что идет дождь, представлять что он идет, или отрицать что он идет. Каждый тип опыта --- восприятие, воображение, желание, воспоминание --- направлен на объект особым образом, и центральная задача феноменологии состоит в анализе этих различий и их взаимосвязей.

\section{Восприятие и образы}

В своих классических работах Гуссерль подчеркивает первичность перцептивной интенциональности. Например, мы могли бы обсудить, какой торт купить на день рождения: ''Отелло'' или ''Захер''? Об этих тортах мы можем думать, говорить, осознавать их даже при их в непосредственной пространственной близости отсутствии. Можем также рассмотреть фотографии торта, чтобы узнать больше о его составе и внешнем виде, при этом будучи направлены на тот же торт что и до этого, но теперь обладающий некоторым графическим присутствием. Наконец, торт мы можем купить или испечь, увидеть и попробовать. Во всех этих трех случаях мы направлены на один и тот же вид торта / вид торта, но проявляет он себя по-разному. Интенциональные акты можно ранжировать в соответствии с их способностью дать объект наиболее непосредственно. Рассмотрим, например, языковые акты: когда мы просто думаем и говорим о торте, он, конечно, объект нашего внимания, но не дан интуитивно. Изобразительные же акты уже имеют некоторое интуитивное содержание, но, подобно языковым, лишь косвенно направленны на объект: языковые акты направлены через конвенциональные представления --- языковые знаки, изобразительные --- через изображение, имеющее определенное сходство с объектом. И только акты восприятия дают нам объект непосредственно, именно этот тип направленности, по Гуссерлю, предоставляет нам сам объект в его телесном присутствии.

Утверждается, впрочем, не только представленность объекта перцептивной направленностью в привилегированном виде, но и более фундаментальный статус такой интенциональности в сравнении с интенциональностями сложными --- через воспоминание или изображение. Эпизодические воспоминания подразумевают воспоминание предыдущего опыта восприятия, воспоминание основано на восприятии. С изображениями сложнее. Стоя в Музее изящных искусств в Вене перед портретом императора Максимилиана работы Дюрера, на какой объект я направлен? Возможно, на изображение-вещь, т.е. обрамленный холст со слоями краски на нем. Возможно, на изображение-объект, т.е. на изобразительную репрезентацию и ее эстетические качества. Или, возможно, на изображение-субъект, т.е. на изображенное --- обыкновенно именно оно привлекает наше внимание первым. Может очаровать выразительный профиль Максимилиана, его роскошное платье или гранат, который он держит в левой руке --- кажется, такая работа внимания банальна и очевидна, но она включает довольно сложную стратифицированную форму интенциональности. Что значит быть изображением чего-то [другого], как один объект может изображать другой? Между знаком и обозначаемым нет внутренней связи, но предполагается наличие у изображения некоторых внутренних качеств, ограничивающих изображаемое им. При правильном сочетании буквы с, т, у, л обозначают стул, но это условность и возможно обозначение таким содержанием другого. В английском соответствующими буквами будут c, a, i, r. Но нельзя так просто решить, что на портрете Дюрера изображена банка томатного супа ''Хайнц''. Впрочем, не только наличие определенного сходства делает один предмет изображением другого: две травинки могут выглядеть совершенно одинаково, но это не делает одну из них изображением другой. Более того, сходство --- отношение двунаправленное, тогда как репрезентация однонаправлена: портрет --- изображение Максимилиана, но Максимилиан не изображение портрета. Согласно анализу изобразительной интенциональности Гуссерля, один и тот же физический объект может функционировать как изображение чего-то другого только будучи воспринимаем определенным образом. Формулируя несколько парадоксально: чтобы увидеть изображенное на картине, нужно и видеть и не видеть физически находящееся перед нами. Нужно воспринимать раму и холст, но также и выйти за пределы физически присутствующего чтобы позволить изображенному проявиться. Рамка своего рода окно в изображенное. Например, читая газету и натыкаясь на фотографию Макрона, вы фокусируете внимание на нем, но достаточно нечеткая фотография будет перетягивать направленность на себя, и даже в наилучшем случае интенциональность останется опосредованной и периферийное осознание изображения-вещи останется, иначе это будет иллюзия.

\section{Представление и причинность}

Интенциональность не просто, но многогранна и многослойна. Сознание постоянно занято объектами и событиями от него отличными, будь то через перцептивную интенциональность или через воспоминание о посещении галереи с изображением X с помощью Y. Но как это возможно? Феноменология отвергает субъективность и отсутствие влияния переживаний на внешний мир. Иначе для объяснения интенциональности был бы нужен некий интерфейс между разумом и миром: торт, на который я смотрю, осознается мной, но он, конечно, не присутствует в сознании буквально, как физический объект, а воздействует на мой сенсорный аппарат, и в результате в сознании возникает мысленное представление о торте. Т.е. восприятие предполагает две сущности: внешний сознанию объект и внутреннюю сознанию репрезентацию, и доступ к внешнему опосредован внутренним.

Такая конструкция отвергается всеми феноменологами.

Одна из их критик направлена на введение репрезентативного опосредования, вторая --- на попытку объяснить интенциональность через причинность.

Репрезентационализм, как известно, порождает скептицизм: почему это внутреннее осознание одной вещи должно способствовать осознанию совершенно другой --- внешней, и как вообще определить, представляет ли внутреннее что-то внешнее? Мы же, в конце концов, не можем выйти за пределы своего сознания и сравнить внутреннее с внешним. Кроме того, исследуя перцептивный опыт, следует понимать его как сталкивающий не с изображениями или знаками объектов (за исключением, конечно, восприятия картин или фотографий), а с самими объектами. Видя фотографию Эмпайр-Стейт-Билдинг или портрет императора Максимилиана, мы имеем дело со сложной формой интенциональности: наше осознание одного --- фотографии или картины --- позволяет осознать другое --- здание или человека, т.е. мы направлены не на воспринимаемое, а через него на что-то другое. В обычном же восприятии воспринимаемое --- не знак или изображение другого. Согласно Гуссерлю, прежде функционирования знаком или изображением объект должен быть воспринят, а значит репрезентативная теория восприятия неверна, потому что предполагает то, что пытается объяснить.

Что же касается причинности, не лучше ли объяснить перцептивную интенциональность как форму причинной связи между объектом и субъектом восприятия? Может, причинность --- клей, соединяющий разум и мир, так что сознательное состояние представляет и направлено на объект тогда и только тогда, когда связано с ним причинно? Такой довольно грубый подход сталкивается с очевидными трудностями. Во-первых, глядя в бинокль на далекий холм, я воспринимаю холм. Хотя холм (или свет, отраженный от него) мог повлиять на мою зрительную систему, это влияние было бы не единственной причиной. Почему я не воспринимаю (репрезентирую) линзы бинокля или проксимальное раздражение глазной сетчатки? Во-вторых, понятие причинности слишком грубо для схватывания аспектной природы интенциональности. Входя в комнату и видя потертый коричневый чемодан, я вижу его определенным образом, не весь сразу, в определенном освещении и на определенном фоне. Можно ли объяснить это причинным влиянием чемодана на мою зрительную систему? Следует помнить, что чемодан появляется всегда в определенном контексте и с определенным значением: в зависимости от предшествующего опыта и текущих интересов, один и тот же чемодан может проявлять себя как дорожное снаряжение, как контейнер для старых писем, как пример имеющего обратную сторону объекта, как сувенир или как символ нынешнего кризиса беженцев, и совершенно неясно, как причинное воздействие чемодана на меня способно охватить все эти различные его проявления. Но основное возражение против сведения интенциональности к причинности следующее: реально существующие пространственные объекты в моем непосредственном физическом окружении, могущие причинно на меня воздействовать --- лишь очень малая часть осознаваемого: можно думать об обратной стороне Луны, о квадратных кругах, единорогах, следующем Рождестве или принципе непротиворечия. Как эти отсутствующие, невозможные, фиктивные, будущие или идеальные объекты могли бы воздействовать на меня причинно? Возможность мышления о несуществующем кажется решающим аргументом против моего осознания объекта только лишь через его причинное на меня влияние.

Феноменологи предлагают говорить о направленности на реально существующие объекты: вкус лимона, запах кофе, холод кубика льда --- качества наличных объектов, а не каких-то ложных психических их коррелятов. Не следует говорить о перцептивном опыте как о внутреннем киноэкране, объекты не должны внутренне воспроизводиться или репрезентироваться для переживания --- перцептивный опыт дает нам объект сразу и непосредственно. Мы не переживаем репрезентации, перцептивная интенциональность не форма репрезентации --- она форма презентации.

Следует также отказаться от попытки объяснения отношения разума и мира в причинных терминах. Объекты в мире взаимодействуют причинно, но разум --- не просто объект, и его с миром отношения нельзя приравнять или ассимилировать с причинными отношениями. Подход к разуму как к такому же объекту как, например, торт, мешает раскрыть и прояснить некоторые из наиболее интересных аспектов сознания, включая его истинное эпистемическое и онтологическое значение.

Итак, феноменологи интересуются анализом структуры интенциональности для прояснения отношений между разумом и миром. Они не слишком заинтересованы отношениями между разумом и мозгом, например, и не согласны с редуктивным объяснением интенциональности через неинтенциональное. Они утверждают, что сознание характеризуется внутренней несводимой интенциональностью и что фундаментальную роль играет значение, а не причинность: камень может удариться об окно и тем самым привести к трещине в нем, эти объекты причинно связаны, но от этого не становятся одним целым. И даже если торт причинно влияет на меня, мое его перцептивное осознание основано не на причинности, а на его для меня значении.

\section{Разум и мир}

Об интенциональности можно рассказать гораздо больше, но для наших целей достаточно указания на интенциональную природу сознания и подчеркивание этим его самопревосходящего характера: разум не замкнут в ожидании причинного воздействия извне для связи с миром, и ошибочно считать мир --- чем-то внешним, а сознание --- расположенным внутри некой сферы с внутренними представлениями мира. Сознанию не нужно буквально выходить за пределы себя для достижения мира, так же как и не нужно ему этот мир как-то поглотить и переварить. Сознание не вместилище и не особое место, сознание определяется открытостью. Мы не внутри изолированной субъективной сферы, из которой время от времени выходим в мир --- интенциональная открытость сознания является его неотъемлимой частью. Это хорошо выражено в знаменитом отрывке из книги Хайдеггера ''Бытие и время'':

\smallskip
\leftskip1cm\relax
\rightskip1cm\relax

Будучи на нечто направлен и его схватывая, Dasein не заключен в некой сфере и выходит за пределы, но всегда уже ''вне'' вместе с этим обнаруженным в мире нечтом. Dasein не покидает никакую внутреннюю сферу, пребывая с познаваемым и определяя его характер. Даже в этом ''бытии вне'' вместе с объектом Dasein в некотором смысле ''внутри'', существует как знающее бытие-в-мире.

\leftskip0cm\relax
\rightskip0cm\relax
\smallskip

Хайдеггер не говорит об интенциональном субъекте, о субъективности или сознании, а предпочитает использовать термин ''Dasein'', состоящий из ''Da'' --- ''там'' и ''sein'' --- ''бытие''. Этот терминологический выбор частично мотивирован желанием избежать некоторых вводящих в заблуждение коннотаций, присущих традиционным терминам, и частично --- желанием подчеркнуть необходимость нахождения самого нашего существа в мире, необходимость его с ним связи. Феноменологическое исследование показывает, что традиционные категории, такие как внутреннее и внешнее, не могут адекватно отразить отношения между Dasein и миром. Поскольку Dasein всегда уже среди вещей, то не имеет внешнего и бессмысленно тогда говорить о наличии у него внутреннего. Схожие идеи можно найти и у Гуссерля, утверждающего, что поверхностное разделение между внутренним и внешним не подходит для корректного понимания интенциональности и что субъект ни в мире, ни вне него, и так же мир ни в субъекте, ни вне него. Феноменология противится самому разделению между внутренним и внешним, как выразился Мерло-Понти: ''Мир полностью внутри меня, а я полностью вне себя''.

Возвращаясь к изобразительной интенциональности, следует заметить не только ее отличие от перцептивной, но и как кажущееся естественным и простым --- видение изображения Эйфелевой башни --- в действительности довольно сложно. Мы обычно воспринимаем объекты как нечто само собой разумеющееся, не обращая внимания вклад нашей интенциональности. Но изображения существуют непростым образом: без разума в мире их бы не было, потому что предметы становятся изображениями только благодаря особой интенциональности. Радикальное феноменологическое утверждение состоит в принятии этого и в более широком смысле. Как пишет Гуссерль в одном из ранних текстов:

\smallskip
\leftskip1cm\relax
\rightskip1cm\relax

Сознаваемые объекты не находятся в сознании как в ящике в ожидании когда их там найдут и схватят; [...] они конституируются как то, чем являются для нас и как то, чем они считаются для нас в различных формах объектной интенциональности.

\leftskip0cm\relax
\rightskip0cm\relax
\smallskip

В текстах феноменологов термин ''конституция'' носит технический характер. Конституция не значит создание: сознание не создает конституируемые им объекты. Они и не выводятся из его действия: составленность воды из водорода и кислорода, а не, например, гелия и ксенона, нельзя объяснить сознанием. Конституирование не значит формирование по своему образу, это допускающий манифестацию или проявление объектов и их значения, процесс, и в процессе этом, конечно, значительно задействованно сознание.

Мы обыкновенно поглощены событиями в мире и, сосредотачиваясь на объекте, но не на том, как он проявляется, не останавливаемся и не задаемся вопросом, как вещи могут являться нам такими и с таким значением. Но чтобы прояснить понятия воспринимаемого объекта, запоминаемого события, оцениваемого положения вещей, нельзя игнорировать интенциональные состояния --- восприятие, припоминание и суждение. Задачей феноменологии с самого начала был отказ от наивности повседневного понимания и исследование корреляции между действием и объектом, между \textit{cogito} и \textit{cogitatum}.

Следует различать интенциональный опыт и интенциональный объект. Они, впрочем, по существу связаны. Корректное феноменологическое исследование интенциональности должно не просто учитывать субъективную сторону --- перцептивный или эмоциональный опыт, например --- но также и исследовать их объективные корреляты --- воспринимаемый или желаемый объект. То же верно и в обратном порядке: нельзя должным образом проанализировать интенциональный объект, не принимая во внимание субъективный коррелят --- интенциональный акт. Основной задачей феноменологии Гуссерль ставил именно объяснение этого отношения между субъективным и объективным: величайшие и наиболее важные проблемы феноменологии связаны со способами конституирования сознанием объективностей разной степени научности, ''поэтому речь идет об исследовании в наиболее всеобъемлющей универсальности конституирования свойственным сознанию способом объективных единств всякой области и категории''.

Интенциональность интересна феноменологам не из-за субъективного опыта как такового, а потому как для понимания статуса физических объектов, математических моделей, химических процессов, социальных отношений, культурных продуктов и т.д. необходимо выяснить, как они могут являться так и с таким значением, и для этого, конечно, нужно учитывать субъекта, которому они являются. Тщательное философское исследование воспринимаемых, оцениваемых, рассматриваемых объектов приводит нас к эмпирическим структурам, с которыми коррелируют способы их проявления: к актам представления, восприятия, суждения и рассмотрения, т.е. к субъекту / субъектам, в отношении которого / которых объект и должен в своем проявлении пониматься.

Феноменологическая позиция предполагает тематизацию данности объектов, но фокус здесь не только на объектах как они даны, а также на структуре соответствующего объектного опыта, что позволяет осознавать интенциональность, благодаря в том числе которой объекты даны именно так. Исследуя проявления объектов, мы говорим о себе как о тех, кому они являются. Тема феноменологического анализа, значит, не обходится без мира --- феноменология не игнорирует мир, и ее интерес к сознанию обусловлен его раскрытием мира: чтобы понять, как мир может быть дан так, как он дан, с таким значением, феноменология приходит к исследованию раскрывающего действия интенционального сознания. Исследование интенциональности не то же, что исследование психологии, и призвано дать правильное понимание реальности и объективности.

Одна из причин центрального места теории интенциональности в проекте Гуссерля --- способность изучения миронаправленности сознания дать понимание не только структуры субъективности, но и природы объективности, потому что возможность чего-то вроде сознательного обретения мира говорит не только о сознании, но и о мире. Конечно, такой способ говорить о сознании --- как о делающим возможным всякое проявление мира конститутивном измерении, как о ''месте'', в котором он может проявиться и выразить себя --- такой способ кардинально отличен от попыток рассмотрения его с научной точки зрения как просто еще одного психического или физического объекта в мире.

Настаивая на необходимости одновременного исследования сознания и мира, феноменология предлагает колеблющую и даже подрывающую традиционно проводимое различие между эпистемологией и онтологией точку зрения, ведь обыкновенно вопрос касательно понимания и познания мир отделяется от вопросов касательно его природы, реальности. Соблазнительно простой шаг здесь состоит в указании на все же необходимость ''вычесть'' всякий вносимый нами субъективный вклад для ответа на второй вопрос --- необходимость исключения субъективного для получения взгляда на реальность ''из ниоткуда''. Но сосредотачиваясь на явлениях, феноменология одновременно анализирует и наш способ понимания и восприятия мира, и сами объекты и способы их проявления. Поэтому в ''Бытии и времени'' Хайдеггер утверждает возможность онтологии только как феноменологии и ключевую роль анализа нашего бытия-в-мире для всякого онтологического исследования.

Феноменологи оспаривают лишь случайность связи между миром и субъективностью, их подобие двум деталькам LEGO, которые можно соединить или разъединить. Интенциональность значит сущностную открытость разума и сущностную манифестируемость мира: чтобы нечто считалось реальным, оно должно быть чем-то, с чем мы можем иметь дело, с чем мы можем столкнуться, хотя, конечно, сам способ встречи с ним может варьироваться: это может быть перцептивное знакомство, практическое участие или научное исследование, или другие возможные формы. Утверждение же, что луна, нейрон, колода карт или общественный ритуал обладают непостижимой и скрытой истинной сущностью, что в действительности они нечто совершенно отдельное от всякого практического контекста, сети значений, теоретической структуры, и что никакая, значит, теоретическая точка зрения не достигает своей цели --- утверждение это не только сбивает с толку, но и попросту эпистемологически наивно. На каком основании и с какой точки зрения оно может быть оправдано? Мы же не можем искоса взглянуть  на свой опыт и увидеть, в какой степени он соответствует реальности, и дело не в трудности достижения такой точки зрения, а в бессмысленности самой идеи ее возможности. Любое понимание реальности по определению перспективно, так что будто бы стирание нашей точки зрения не приближает к миру, а напротив --- мешает что-либо вообще понять о нем.

\section{Реализм и идеализм}

Настаивая на взаимозависимости и нераздельности разума и мира, феноменологи не сводят его к внутриментальным модификациям или конструкциям, а сочетают отличность мира от разума с его с разумом взаимной тесной связью. В одном из ранних текстов Сартра можно найти яркую формулировку этой мысли:

\smallskip
\leftskip1cm\relax
\rightskip1cm\relax

Гуссерль настойчиво утверждал невозможность растворения вещи в сознании: вы, конечно, видите, например, это дерево, но видите вы его там, где оно есть --- на обочине дороги, среди пыли, в восьми милях от побережья Средиземного моря. Оно не может войти в сознание, ведь не имеет той же природы. [...] Но Гуссерль не реалист: дерево это не абсолют, вступающий с вами в контакт. Сознание и мир даны сразу: будучи внешним сознанию, мир по существу соотнесен с ним.

\leftskip0cm\relax
\rightskip0cm\relax
\smallskip

Научная литература полна дискуссий об идеалистичности феноменологии от такого взгляда. Левинас однажды отметил, что идеализм Гуссерля не о замкнутом в себе и знающем лишь свои состояния субъекте, а о его интенциональности и потому всему открытости. Многое зависит от используемых определений: определяя идеализм как знание субъектом лишь собственных состояний, и реализм --- как непосредственное осознание объектов с типом бытия совершенно иным чем разум, в мире, можно говорить о Гуссерле как о скорее реалисте. Но, как ранее уже отмечалось, феноменологи совершенно точно выступают против объективизма, против понимания реальности как совершенно независимой и ее нами лишь самонадеянного отражения.

Так считал не только Гуссерль: Хайдеггер и Мерло-Понти отрицали самодостаточную природу разума и утверждали его по своей сути включенность в мир. Они, впрочем, защищали также и связь мира с разумом, т.е. говорили об отношении между разумом и миром как о внутренним --- конституирующим свои реляты, а не внешним --- основанном на причинности. В курсе лекций ''Основные проблемы феноменологии'' 1927 года Хайдеггер пишет:

\smallskip
\leftskip1cm\relax
\rightskip1cm\relax

Мир существует, есть, только если существует, есть Dasein: только если есть мир, если есть присутствие как бытие-в-мире, есть понимание бытия, только если это понимание существует --- только тогда внутримировое раскрывается как существующее. Миропонимание как Dasein-понимание есть самопонимание, я и мир принадлежат друг другу в единой сущности --- Dasein, они не два как субъект и объект или как я и ты, но суть основное определение самого присутствия в единстве структуры бытия-в-мире.

\leftskip0cm\relax
\rightskip0cm\relax
\smallskip

Схожая приверженность взаимозависимости разума и мира присуща и Мерло-Понти, который в конце ''Феноменологии восприятия'' заявляет:

\smallskip
\leftskip1cm\relax
\rightskip1cm\relax

Мир неотделим от субъекта, но от субъекта который не что иное как проект мира, и субъект этот неотделим от мира, но от мира который сам и проецирует''. Субъект есть бытие-в-мире, и мир остается ''субъективным'', поскольку текстура его и артикуляции очерчиваются движением трансцендентности субъекта.

\leftskip0cm\relax
\rightskip0cm\relax
\smallskip

В центре внимания феноменологии находится пересечение разума и мира, и ни один из них не может быть понят отдельно от другого. Мы то что мы есть в зависимости от нашей вовлеченности в мир, и мир как фундаментальный смысловой контекст таков каков есть только благодаря нашей в него вовлеченности, и вопрошать об одном без другого все равно что вопрошать что есть фон сам по себе, независимо от переднего плана.

\chapter{Методологические соображения}

Один из наиболее противоречивых и обсуждаемых вопросов феноменологии касается метода. Каков феноменологический метод, если, конечно, он вообще существует? Вопрос спорный не только из-за широко распространенных разногласий и весьма разрозненных ответов, но и потому что это один из наиболее неверно понимаемых аспектов феноменологии и ее критиками, и ее сторонниками.

\section{Эпохе и редукция}

Гуссерль известен утверждением необходимости, для выполнения феноменологией своих задач, определенных методологических шагов. Он настаивал даже, что игнорирование этих шагов значит невозможность понять, что такое феноменология. О каких же шагах речь? Давайте остановимся на двух хорошо известных. В ''Логических исследованиях'' Гуссерль заявляет: ''Совершенно точно нельзя довольствоваться лишь словами [...]. Смыслов, внушенных только отдаленными, смутными, недостоверными интуициями --- если интуициями вообще --- недостаточно: следует вернуться к самим вещам''. Но что именно следует вынести за скобки или отстранить, и как такая процедура позволит вернуться к самим вещам? Здесь мнения расходятся. Согласно одной интерпретации, возвращение к самим вещами значит отход от теорий, интерпретаций и построений: нам следует вынести за скобки наши предвзятости, привычки мышления, предрассудки и теоретические предположения, и вместо вытаскивания на сцену всего этого теоретического багажа, осуществить непредвзятый поворот к объектам, позволив им проявить себя такими, какие они есть. Следует, значит, сосредоточиться на вещах как они встречаются в опыте, а не на том, какими мы их представляем, и уже затем переходить к аккуратным описаниям и определениям. В таком прочтении феноменология скорее описательна, чем дедуктивна или спекулятивна, и ее ядром является строгий интуитивный метод. Иногда описание понимается как попытка игнорировать особенности объекта и выделение его только существенных особенностей, иногда же утверждается необходимость такого метода быть как можно более подробным для учета и схватывания уникальности конкретного исследуемого явления. Одно из ранних (дофеноменологических) проявлений этого стремления можно найти у Флобера:

\smallskip
\leftskip1cm\relax
\rightskip1cm\relax

Мы привыкли всякий раз используя глаза вспоминать, что другие думают о визуально нам явленном. Даже в малейшей вещи есть неизведанное, и мы должны его найти. Для описания пылающего огня или дерева на равнине мы должны оставаться перед этим огнем или деревом до тех пор, пока они не перестанут напоминать какой-то другой огонь или дерево.

\leftskip0cm\relax
\rightskip0cm\relax
\smallskip

Эпохе можно понимать и иначе: следует вынести за скобки или игнорировать не только традиционные теории и предрассудки, но также --- что более важно --- нашу привычную, естественную озабоченность объектами и событиями в мире. Цель феноменологии --- выявление обыкновенно упускаемых из виду из-за сосредоточенности на объектах, аспектов и измерений нашей субъективной жизни. Т.е. феноменология стремиться расширить сферу нашего внимания для тематизации и описания до сих пор незамеченных аспектов внутреннего опыта, и потому и критики, и сторонники феноменологии приходят к выводу о ее сходстве с интроспективной психологией.

Из этих двух интерпретаций ошибочны обе. Не то что бы они вовсе не содержат истины, но упускают важную часть: феноменология не поворот к объекту, и не возвращение к субъекту, и ошибочно полагать главной ее целью предоставление детальных описаний. Приведенные интерпретации не только неспособны охватить рамки феноменологического анализа --- цель которого, как уже подчеркивалось, не в исследовании отдельно объекта или отдельно субъекта, отдельно мира или отдельно разума, а исследование их взаимосвязи или корреляции --- они также не могут признать собственно философскую природу феноменологии. Не случайно Гуссерль и Шелер отвергли чисто описательную, лишенную систематических амбиций ''феноменологию книжки с картинками''. Одно только накопление различных описаний, будь то индивидуальный опыт --- то что чувствую здесь и сейчас --- или отдельные объектов, или их более инвариантных сущностных структур --- не заменяет систематической и аргументированной работы, которую следует выполнять феноменологам. Утверждать достоинством феноменологии мотивацию внимания к тому, как выглядят вещи --- восход солнца или аромат кофе --- значит упустить совершенно фундаментальным образом наиболее важную ее часть.

Прежде чем сформулировать правильную интерпретацию, давайте кратко рассмотрим еще одно неверное, но крайне влиятельное толкование, согласно которому эпохе значит методологическую или метафилософскую, а не метафизическую природу феноменологии. Идея в том, что феноменологам следует взять за скобки и исключить из рассмотрения реально существующий мир, сфокусироваться на явлениях, на том, как вещи выглядят и какое значение для нас имеют, что следует фокусироваться, конечно, не только на опыте, но на объектах лишь постольку поскольку они в опыте, и не важно, существуют ли они в действительности.

Согласно этой интерпретации, цель феноменологического ''заключения в скобки'' --- ограничение рамок исследования: определенные вопросы, которым феноменологам не положено заниматься, просто исключаются из рассмотрения. Так, мы могли бы верить, что направлены на что-то эстраментальное, на что-то трансцендентное, не содержащееся в сознании, но будучи феноменологами, должны исследовать это убеждение и наш опыт объектов, но не вправе говорить об их бытии. Т.е. как феноменологу, мне следует утверждать, что я ощущаю лимон, что он мне является, что мне кажется что передо мной лимон, но не следует говорить о его в действительности существовании, ведь это значило бы переход от феноменологии к метафизике.

Эта интерпретация неверна и очень влиятельна. И Хайдеггер, и Мерло-Понти совершенно ясно заявляют о своих онтологических приверженностях и о степени приверженности их феноменологической работы к бытию. Как заявляет Хайдеггер, ''рядом с феноменологией нет онтологии, но научная онтология есть ничто иное как феноменология''. Можно, конечно, утверждать, что феноменология после Гуссерля претерпела радикальную трансформацию. Его концепция феноменологии действительно включала вышеупомянутые ограничения, все более поздние феноменологи отказались от его методологических приемов и радикально изменили проект феноменологии. Было написано множество книг о спорных отношениях между гуссерлианской и постгуссерлианской экзистенциальной или герменевтической феноменологией и здесь не место рассматривать эти отношения детально, но позвольте открыто высказать свое мнение по этому вопросу. Я на стороне защитников тезиса и непрерывности и считаю, что есть нечто вроде феноменологической традиции с набором общих тем и проблем, объединявших и объединяющих ее сторонников.

Наиболее прямой способ защиты этой интерпретации --- показать, что Гуссерль все же интересовался миром и истинным бытием, и что цель эпохе не в исключении их из рассмотрения. Правильный способ интерпретировать эпохе --- рассматривать его как подразумевающее не исключение реальности, а как приостановку определенного догматического к ней отношения, действующего не только в позитивных науках, но также и пронизывающего нашу повседневную дотеоретическую жизнь. В чем заключается эта естественная установка? Например, в простом принятии как должного существования мира также и независимо вне нас. Сколь бы естественно и очевидно не было думать о реальности как о самодостаточной, философия, представляя радикальную форму критического объяснения, не может принимать этот вид реализма как нечто само собой разумеющееся. В стремлении заслужить репутацию формы радикального вопрошания, философия не может предполагать ответ --- напротив, желая принять феноменологическую позицию, следует сделать шаг назад от нашего наивного и неисследованного погружения в мир и приостановить нашу автоматическую веру в независимое от разума существование мира. Приостановив же эту установку и тематизировав раскрытие и исследование реальности всегда с той или иной точки зрения, мы не только не потеряем ее из виду, но и напротив --- впервые сделаем ее доступной философскому исследованию.

Но почему же тогда многие интерпретировали эпохе как воздержание от онтологических претензий, как игнорирование касающихся бытия вопросов? В некоторых неверных интерпретациях отчасти виноват сам Гуссерль. В разделе 52 ''Кризис'', он сначала утверждает ограничение эпохе всякого интереса к миру:

\smallskip
\leftskip1cm\relax
\rightskip1cm\relax

Запретен всякий интерес к бытию, действительности или небытию мира, то есть всякий теоретически ориентированный на познание мира и даже практический в обычном смысле с его зависимостью от предпосылок его ситуационных истин.

\leftskip0cm\relax
\rightskip0cm\relax
\smallskip

Но  страницей позже поясняет, что это первоначальное представление вводит в заблуждение:

\smallskip
\leftskip1cm\relax
\rightskip1cm\relax

При соответствующей эпохе переориентации ничего мирового не теряется --- ни интересы, ни цели, и значит, цели познания также остаются. Но для всех них выявляются их существенные субъективные корреляты, и так раскрывается полный онтический смысл объективного бытия и, значит, всей объективной истины.

\leftskip0cm\relax
\rightskip0cm\relax
\smallskip

Далее он добавляет, что одно из распространенных заблуждений об эпохе --- полагание за ним ''отказа'' от ''всяких естественных человеческих жизненных интересов''. Так называемое исключение мира --- в действительности исключение определенного наивного предрассудка касательно метафизического статуса этого мира: феноменологу следует лишь прекратить наивно постулировать мир. Эпохе не предполагает ''приостановки'' мира, а направлено на открытие фундаментальной корреляции между разумом и миром. И Гуссерль весьма настойчиво подчеркивает необходимость расширять, а не ограничивать, область исследований, даже уподобляя в ''Кризисе'' действие эпохе переходу от двухмерной к трехмерной жизни: выполняя эпохе, феноменолог должен совершить своего рода гештальт-сдвиг, изменение перспективы, расширяющее понимание.

Впрочем, эпохе --- лишь первый шаг. ''Вынося за скобки'' и приостанавливая естественную установку и так перестав принимать реальность за бесспорную отправную точку, феноменолог обращается к данному и способам которыми оно дано. Но анализируя это, он приходит также к открытию интенциональных действий и эмпирических структур, только в отношении которых всякий являющийся объект и может пониматься. Он признает свой собственный, субъективный, вклад и интенциональность, благодаря которой объекты являются именно так и имеют именно такой смысл. Говоря о трансцендентальной редукции, Гуссерль имеет в виду систематический анализ этой корреляции между субъективностью и миром, и анализ этот более продолжительный и ведет от природной сферы обратно к ее трансцендентальным основам. И эпохе, и редукция, значит, могут рассматриваться как элементы философской рефлексии, ставящей целью освобождение от естественного (натуралистического) догматизма и осознание нашей собственной конститутивной вовлеченности.

\section{Трансцендентальная философия}

Настаивая на необходимости такого рода рефлексивного движения, Гуссерль подтверждает свою принадлежность к традиции трансцендентальной философии. Цель этого философского подхода не в каком-либо прямолинейном исследовании объектов в мире для раскрытия новых фактов, как это делается ботаниками или морскими биологами --- правильное философское исследование реальности состоит не в такой инвентаризации, а в выяснении условий реальности чего бы то ни было. Вместо наивного предположения объективного характера мира, следует сначала объяснить, как возможно объективное, как мир возможен объективным и данным в опыте трансцендирующим этот опыт.

Принимая феноменологическую установку, мы не отворачиваемся от мира внутрь себя, а продолжаем интересоваться объектами мира, но уже не наивно, а лишь так как он дан, как коррелятом опыта. Интерес Гуссерля к сознанию обусловлен его раскрытием мира: оно в мире не только объект, но и субъект и, значит, необходимое условие возможности проявления всякой сущности.

Феноменология, таким образом, направлена на не ограниченное исследование психологической области, а на правильное понимание реальности и объективности, и говорить о феноменологии как об оставляющей бытие и реальность прочим дисциплинам --- значит не уважать собственные воззрения Гуссерля по этому вопросу. Как он заявляет в разделе 23 ''Картезианских медитаций'' (1931), темы существования и несуществования, бытия и небытия --- всеобъемлющи для феноменологии.

Эпохе и редукция --- это не попытка заставить реальность исчезнуть из поля зрения, а напротив --- попытка философски ее исследовать. В своем раннем курсе лекций ''История понятия времени: пролигомены'' Хайдеггер рассматривает феноменологическую методологию Гуссерля и предлагает следующую весьма точную характеристику:

\smallskip
\leftskip1cm\relax
\rightskip1cm\relax

Такое ''заключение в скобки'' сущности ничего от нее не умаляет и не предполагает ее отрицание: скорее это изменение перспективы делает бытие сущности видимым. Такая феноменологическая приостановка трансцендентального тезиса нужна лишь для одного --- сделать сущность присутствующей в отношении ее бытия. ''Приостановка'', таким образом, понимается неправильно, покуда видится разрыванием связи рефлексии с сущностью, ведь на самом деле именно для выяснения бытия данной сущности такая приостановка и делается.

\leftskip0cm\relax
\rightskip0cm\relax
\smallskip

Гуссерль неоднократно настаивал на невозможности понять феноменологию вне эпохе и редукции. Хайдеггер и Мерло-Понти редко упоминали эти методы, но потому ли, что отвергли методологию Гуссерля, или может потому, что приняли ее как нечто само собой разумеющееся? В разделах 27-33 ''Идей I'' Гуссерль подробно описывает естественную установку и утверждает, что фундаментальные структуры нашего отношения к миру и особый характер нашей собственной субъективности буду оставаться скрытыми, покуда мы не откажемся от наивной, дофилософской, позиции, и только приостановив естественную установку, мы обнаружим, что субъективность --- нечто большее, чем бытие объектом в мире.

Перейдя от ''Идей I'' к ''Бытию и времени'', можно обнаружить несколько схожую манеру рассуждения у Хайдеггера: для него повседневное существование характеризуется самозабвением и самообъективацией. Мы склонны позволять нашему самопониманию роководствоваться, и, следовательно, прикрываться, здравым пониманием происходящего в мире. Феноменология значит борьбу против такого нивелирующего самопонимания, и потому Хайдеггер в ''Бытии и времени'' пишет, что феноменологическому анализу свойственна определенная жестокость: раскрытие им бытия Dasein можно добиться только в прямой конфронтации с тенденцией Dasein прикрывать вещи, оно должно быть вырвано, захвачено у Dasein. Хайдеггер говорит также о соответствии нашей обыденной жизни общепринятым нормам и стандартам, отчего все будто бы уже понято и истолковано другими, все оценки и суждения воспринимаются некритически, и оставляя это все как есть, мы чувствуем себя в безопасности, вне нужды решения фундаментальных и сложных вопросов. Но некоторые события --- сильная тревога, например --- могут разрушить привычность, сделав даже знакомое и уютное жутким, и тогда уже не получится так легко полагаться на общепринятые интерпретации. Тревога в этом смысле --- событие, нарушающее естественную бездумность и подталкивающее к философским вопросам.

В ''Феноменологии восприятия'' Мерло-Понти пишет:

\smallskip
\leftskip1cm\relax
\rightskip1cm\relax

Будучи совершенно связаны с миром, уловить самих себя мы можем только приостановив это движение, отказавшись от участия в нем, увидеть его вне своего в нем участия. Это не отказ от достоверности здравого смысла и естественной установки --- напротив, это извечное направление философии --- это от их как предпосылок всякой мысли принятия как нечта саморазумеющегося, принятия незаметного и потому необходимого к приостановке, и лишь такое воздержание от них заставит их проявиться. Возможно, лучшая формулировка редукции --- предложенная помощником Гуссерля Ойгеном Финком --- говорившем об изумлении, заинтересованности миром. Рефлексия не уходит от мира к единству сознания как основы мира, а, скорее, отступает назад, чтобы стало видно возникновение трансцендентности и ослабились и оттого проявились связывающие нас с миром интенциональные нити; только так мир сознается и сознается тогда как странный и парадоксальный.

\leftskip0cm\relax
\rightskip0cm\relax
\smallskip

В этой же книге Мерло-Понти будет утверждать необходимость порывания с привычным восприятием мира для правильного его понимания, и невозможность корректного исследования сознания до вопрошания об абсолютном статусе мира.

Ни Хайдеггер, ни Мерло-Понти не согласились бы максимально точным и тщательном описании объектов или опыта как задачи феноменологии --- она не занимается исследованием явлений во всем их фактическом многообразии. Философия для них характеризуется совершенно иной, нежели свойственная позитивным наукам, позицией. Эпохе и редукция --- термины Гуссерля для обозначения рефлексивного движения, необходимого для достижения философского мышления, и несмотря на их разногласия касательно деталей программы Гуссерля, и Хайдеггер, и Мерло-Понти совершенно привержены этому рефлексивному шагу. Позитивная наука, согласно ним, принимает определенные идеи о независимой от разума природе реальности как нечто само собой разумеющееся и, похоже, по умолчанию считает их неуязвимыми для критического анализа. Феноменология ставит под сомнения такой объективизм и предлагает исследовать все объекты, научные открытия, культурные достижения, социальные институты и т.д., сфокусировавшись на способах их проявления. Если феноменологам должно заняться фундаментальными онтологическими вопросами, то, согласно Хайдеггеру, следует продолжить исследование понимания бытия Dasein и исследовать бытие ''постольку, поскольку оно входит в умопостигаемость Dasein''.

Отношения между Гуссерлем и постгуссерлианскими феноменологами противоречивы: расходятся мнения касательно степени долга Хайдеггера и Мерло-Понти перед Гуссерлем. Карман утверждает, что ''фундаментальная онтология Хайдеггера не может быть понята как лишь дополнение или продолжение философии Гуссерля, не говоря уже об этой философии ''переводе'''', Мерло-Понти указывает, что ''Бытие и время'' --- не что иное, как объяснение гуссерлевской идеи о проживаемом мире. И хотя Мерло-Понти сам неоднократно подчеркивал свой долг перед Гуссерлем и время от времени представлял свою работу попыткой раскрытия смысла поздней философии Гуссерля, попыткой осмысления его ''необдуманной мысли'' --- несмотря на это, многочисленные исследователи Мерло-Понти настаивали на лишь экстраполяции его собственной философии под таким вот ''гуссерлевским'' ярлыком.

Между Гуссерлем, Хайдеггером и Мерло-Понти, есть, конечно, существенные различия, и как бы первый не повлиял на остальных, они обязаны также и другим выдающимся деятелям философской традиции, включая Аристотеля, Декарта, Кьеркегора, Ницше, Бергсона и Сартра. Но большая часть разногласий Хайдеггера и Мерло-Понти с Гуссерлем происходит в пределах горизонта общих допущений --- это имманентная критика, внутренняя по отношению к феноменологии, а не порывающая с ней. Иными словами, знакомство с Гуссерлем необходимо для понимания и оценки феноменологического аспекта мышления Хайдеггера и Мерло-Понти.

\chapter{Наука и проживаемый мир}

Для многих феноменологов задача состоит не в описании эмпирических а фактических особенностей, а в исследовании существенных структур, характеризующих наш опыт, их коррелятов и связи между ними. Философов интересуют необходимые и инвариантные свойства, и уже затем, возможно, необязательные и случайные: характеризующее восприятие и желание в целом важнее, чем, скажем, разница между восприятием десятицентовика и пенни или разница между желанием красного вина и белого вина. Восприятие десятицентовой монеты и восприятие пенни в чем-то различны, но имеют и общее, делающее оба их актами восприятия, а не, например, воображения или воспоминания, и на раскрытие именно этих особенностей нацелен философский анализ. Как такое раскрытие возможно?

Прежде ответа позвольте подчеркнуть, что способность отличать случайное от существенного и необходимого занимает центральное место в философии и науке, но используется и в повседневной жизни: стоя в аптеке и глядя на разнообразные модели зубных щеток, я сразу понимаю, что это разные зубные щетки, хотя и отличающиеся по цвету и размеру, и точно так же, ища конкретную книгу, я легко ее узнаю, даже найдя в новом месте или положении, потому что сразу осознаю несущественность ее пространственного положения для ее идентичности.

\section{Эссенциализм}

Один из инструментов, использованных Гуссерлем в поисках существенных структур --- эйдетическая вариация: взятие отправной точкой данного и затем использование воображения для изменения исследуемого и так постепенного устранения его случайных свойств и достижения, значит, свойств необходимых --- не могущих быть измененными без того чтобы объект исследования перестал быть тем, что он есть. У книги, например, может быть другой цвет и дизайн обложки, другое количество страниц, размер, вес и т.д. --- все это известно из опыта, или навеяно воображением. Конечным результатом будет определенный набор свойств, принадлежащих всякой книге, всякую книгу книгой делающих. Эйдетическая вариация помогает раскрыть составляющие сущность объекта инвариантные структуры.

Хотя Гуссерль иногда говорит об этом процессе как о ведущем к своего рода эйдетической или сущностной интуиции (Wesensschau), важно понимать, что он не предполагает доступа к некоему таинственному и безошибочному источнику познания: эйдетическую вариацию можно рассматривать как своего рода концептуальный анализ, управляемый воображением и преследующий совершенно иную эмпирической или экспериментальной работе цель. Это сложный и открытый процесс, результаты которого в большинстве случаев можно изменить, так что нет никаких претензий на непогрешимость. Идеям всегда свойственна некоторая предварительность, презумпция, они всегда открыты для будущих модификаций в свете нового. Гуссерль иногда отличает так называемые точные сущности от морфологических: изучая математику и прочие точные науки, мы можем определять соответствующие объекты с большой точностью, но объектам из гуманитарных и социальных дисциплин свойственна определенность куда меньшая --- они обладают морфологической сущностью, описание и классификация которой лишь приблизительны. Попытки навязать вещам проживаемого мира подобные свойственным, скажем, геометрии, точность и аккуратность, насильственны. Что же касается феноменологии, Гуссерль, конечно же, не полагал ее целью раскрытие предполагаемой сущностной структуры браков, выборов, кларнетов или дубов --- он исследовал гораздо более фундаментальные темы на гораздо более высоком уровне общности. Каковы общие структуры проживаемого мира, интенциональности, воплощения, временности? Что отличает восприятие от воображения? Что характеризует физический объект как таковой и чем он отличается от математической сущности или психологического процесса?

Возможно, будет полезно разграничить приверженность Гуссерля существованию инвариантных и универсальных структур и две альтернативные, но совершенно противоположные позиции. Во-первых, эссенциализм Гуссерля не следует путать с социально-биологическим эссенциализмом, критикуемым многими учеными в расовых, постколониальных и феминистских исследованиях и согласно которому раса, этническая пренадлежность, пол и т.д. --- фиксированные, сущностные и внеисторические определенности. Согласно этому виду эссенциализма, быть отцом, датчанином или евреем --- значит обладать определенными, фиксированными и неизменными свойствами, общими во все времена для всех отцов, датчан и евреев. Гуссерль хорошо осознавал, насколько подобные понятия в действительности подвержены историческим и культурным изменениям.

Во-вторых, точку зрения Гуссерля следует противопоставить релятивизму некоторых социальных конструктивистов. В своей книге ''Кредо конструктивизма'' Линкольн и Губа, например, утверждают, что релятивизм --- основная онтологическая предпосылка конструктивизма, и что объективность --- химера, и изучаемые гуманитарными науками сущности ''на самом деле'' не существуют, поскольку совершенно условны и существуют только в сознании их созерцающих. Упустить из виду решающие различия между этой и утверждаемой Гуссерлем позициями было бы большой ошибкой.

Учитывая распространенность тех или иных форм эссенциализма как в науке, так и в (истории) философии, важно не приравнивать эйдетическую вариацию с трансцендентальной редукцией --- особенно вопиющей ошибкой была бы отдача приоритета первому в ущерб второму.

В предыдущих главах указывалось на стремление феноменологов вернуться к ''самим вещам'', лозунг этот указывает на критичность и недогматичность их исследования, на избегание метафизических спекуляций и руководство действительно данным а не ожидаемым с учетом некоторых теоретических установок. Метод должен диктоваться предметом исследования, а не научным идеалом. Как отмечает Хайдеггер в ''Что такое метафизика'', ошибочно приравнивать научную строгость к математической точности: нельзя полагать всякую невозможную к с математической точностью анализу область менее ценной или менее реальной.

Подобные идеи можно найти и у Гуссерля, который в работе ''Формальная и трансцендентнальная логика'' предостерегает от опасности ослепить себя методологией и идеалами точных наук, понятых составляющими абсолютные нормы истинного и настоящего. Вместо позволения предопределенным теориям права ограничения возможного опыта, следует их сами подчинить объекту опыта, следует позволить явлениям говорить:

\smallskip
\leftskip1cm\relax
\rightskip1cm\relax

Истинный метод следует природе исследуемого, а не предубеждениям.

Необходимо оставить утверждения о наблюдаемом и так постараться его не отрицать и не замалчивать.

\leftskip0cm\relax
\rightskip0cm\relax
\smallskip

Согласно Гуссерлю, ученые вполне могут использовать более точные измерения, чем, скажем, продавцы на рынке, но точность эта имеет свои ограничения. Желая продать килограмм апельсинов, вы не станете измерять вес в микрограммах. Достаточность и уместность зависит от контекста и не определяется в абсолютных терминах.

Но разве эти максимы не тривиальны? Не совсем. С феноменологической точки зрения реальность сложна и состоит из множества различных онтологических областей, будь то царство идеальностей, природы, культуры и т.д., и всякое этих сфер исследование должно учитывать их особенности и использовать соответствующие им методы, и поэтому феноменология открыто критикует множество частично перекрывающихся позиций, включая редукционизм, элиминативизм и натурализм.

\section{Редуцировать или исключить}

Научный редукционизм мотивируется различными методологическими принципами, в том числе так называемой ''бритвой Оккама'', призывающим не предполагать существования большего числа типов объектов без строгой на то надобности. Т.е. если приходится выбирать между несколькими теориями, каждая из которых рассматривает свою (по-видимому) нередуцируемую часть реальности и одной, которая систематически объясняет эти различные части, сводя более сложные их части к менее сложным, то следует выбрать последнюю, и не только потому, что возрастающее единство, систематичность и простота удобны, но и потому, что рассматриваемая редукция должна обладать наибольшей объяснительной силой. Если одни свойства сводимы к другим, то первые можно объяснить вторыми. Классическим примером является объяснение макросвойств --- температуры, прозрачности, растворимости или эластичности --- через микросвойства --- например, через молекулярный состав. Обыкновенно предполагается, значит, что для ответа на вопрос ''что такое x?'' следует переформулировать вопрос так: ''как можно свести x к тому, что можно понять с помощью физики, химии, нейрофизиологии и т.д.?''. Еще одно предположение --- нереальность так не редуцируемого. Формулировку этой точки зрения можно найти у Джерри Фодора: ''Трудно понять [...], как можно быть реалистом в отношении интенциональности, не будучи при этом в той или иной степени редукционистом. [...] Если предметность (aboutness) реальна, то на самом деле она должна быть чем-то другим''. Предполагается, значит, что только редуктивное объяснение, скажем, сознания, может дать правильное понимание его природы и что сознание реально только если такое объяснение возможно.

Элиминативисты разделяют многие из опасений редукционистов: реальным считают только то, что объяснимо методами и принципами естественных наук. Впрочем, в отличие от редукционистов, они утверждают невозможность сведения сознания к, скажем, нейрофизиологическим процессам, из чего делают вывод о его несуществовании. Для элиминативного материалиста вера в существование чего-то вроде ''желания, страха, ощущения, радости, боли и т.д.'' --- не более чем теоретические предположения, составляющие примитивную психологическую теорию, не соответствующую стандартам современной науки и которую потому следует отвергнуть подобно алхимии или френологии. Сознание, таким образом, не реально, а лишь выдумка наравне с единорогами и йети.

Научный натурализм отличается как методологически, так и онтологически. Его методологическая приверженность заключается в правильности тех и только тех процедур и типов обоснования, которые используются в естественных науках. Все подлинные вопросы, значит --- естественнонаучные, а подлинное знание --- объективное, полученное естественнонаучными средствами. Галилей, которого Гуссерль считал ранним олицетворением этой тенденции, писал:

\smallskip
\leftskip1cm\relax
\rightskip1cm\relax

Философия написана в этой великой книге --- Вселенной, постоянно открытой нашему взору. Но книгу не понять, не поняв вначале язык, не научившись читать буквы, из которых она составлена. Написана она на языке математики, знаки ее --- треугольники, круги и другие геометрические фигуры, без которых по-человечески не понять ни слова из нее, без них попытки ее понять будут подобны блужданию в темном лабиринте.

\leftskip0cm\relax
\rightskip0cm\relax
\smallskip

С исторической точки зрения, подобные соображения часто сочетались с онтологическими убеждениями, согласно которым реальность состоит только из принятых или могущих быть принятыми естественными науками сущностей, свойств и структур. Первоначально утверждалась объективность количественно описываемых характеристик --- формы, размера, веса, и субъективность --- лишенность какой-либо независимой от разума реальности --- цвета, вкуса и запаха. Это классическое разделение между первичными и вторичными качествами с годами радикализировалось, в конечном счете приведя к утверждению лишенность объективности всякого являющегося: весь мир опыта --- субъективная конструкция, иллюзорное изображение скрытой физической природы, и покуда наука стремиться раскрыть истинную природу реальности, ей следует выйти за пределы всего феноменологически данного. Мир, в котором мы живем, который мы знаем из опыта --- совершенно отличен от мира науки, но реален именно второй.

В совокупности эти два убеждения сводятся к естественнонаучности всякой истины и необходимости изучения всего существующего --- включая сознание, интенциональность, смысл, рациональность, нормативность, ценности, культуру, история и т.д. --- методами естествознания, для получения единственно онтологически валидных --- естественнонаучных --- фактов.

Очевидно, эти позиции радикальны. Верность лозунга ''редуцировать или исключить'' значит отсутствие реальной научной ценности у многих найденных в социальных и гуманитарных науках объяснений и реальность многих исследуемых ими объектов и явлений: деньги, марки, симфонии, муниципальные выборы, гражданские войны и т.д.. Трудно представить адекватное объяснение гражданской войны в Сирии --- политическое, культурное, социальное и экономическое явление --- с помощью принципов нейрофизиологии, биологии или физики. Но тогда, согласно лозунгу, гражданской войны в Сирии никогда не было, что, конечно абсурдно.

\section{Проживаемый мир}

Размышляя о статусе мира науки, феноменологи часто подчеркивали важность жизненного, проживаемого мира. Это мир, в котором мы живем, который принимаем как должное в повседневной жизни, дотеоретический мир опыта, с которым мы все знакомы и который обычно не подвергаем сомнению. Этот мир, утверждают феноменологи, необходимо реабилитировать --- он подавлен наукой, историческую и систематическую основу которой составляет. Даже самые точные и абстрактные научные теории опираются на донаучные данные о проживаемом мире. В поисках объективного знания наука сделала добродетелью свою способность выходить за пределы телесного, чувственного и прагматического опыта, превосходить его, но часто упускала из виду степень обусловленности этого самим опытом. В конце концов, зачастую физик исследует и научно определяет именно воспринимаемый повседневный объект --- планетарные тела, которые можно наблюдать на небе, вода, которую можно пить, цветок, которым можно восхищаться и т.д. --- именно такой объект естествоиспытатель исследует в стремлении определить его истинную природу как можно точнее и объективнее. Даже в случаях отдаленности исследуемого от повседневной практики, проживаемый мир остается в планировании и постановке экспериментов, в чтении измерительных приборов, при интерпретации, сравнении и обсуждении результатов с другими учеными: хотя научные теории в своей точности и абстракции стремятся вытеснить конкретный и интуитивно данный проживаемый мир, он всегда остается точкой отсчета.

Ошибочно понимать отношения между проживаемым миром и миром науки как статические: наука опирается на жизненный мир, но также и влияет на него, и постепенно ее теоретические открытия поглощаются и интегрируются в него. Штрассер сравнивал жизненный мир с плодородной почвой --- подобно тому как она питает растительность, проживаемый мир питает теоретические знания, и подобно изменению физических и химических свойств почвы от роста растений, изменяется и жизненный мир под воздействием научных теорий, основу которым дал.

Указание на значимость проживаемого мира не значит критику науки. Феноменология не ставит под сомнение ценность науки, не отрицает способность научных исследований приводить к новым открытиям и расширять наше понимание реальности. Отвергается лишь возможность естественных наук дать исчерпывающее описание реальности. Опять же, речь не о противопоставлении феноменологии количественным методам и исследованиям --- они превосходны в решении количественных вопросов. Для феноменологии вопрос о реальности нечта не зависит от возможности уместить его в прокрустово ложе количественной науки: мир опыта имеет свои собственные критерии достоверности и истины и не обязан ждать одобрения науки. В конце концов, результаты, к которым приходит наука, и повседневный опыт не обязаны противоречить друг другу --- все они могут быть истинны в соответствии со своими собственными стандартами. Для феноменологов различие между миром опыта и миром науки вообще различие не между миром для нас и миром в себе, а между двумя способами проявления одного и того же мира. Мир науки --- не отдельный, автономный мир где-то позади или в глубине проявленного мира. Изучаемый наукой мир тот же, что и мир повседневного опыта --- проявленная реальность, разве что исследуется он в научных терминах. Предстающий перед нами мир --- будь то в восприятии, в повседневных заботах или в научном анализе --- один и тот же, единственный реальный мир, и утверждать существование помимо него некоего закулисного, превосходящего всякую видимость, всякое, значит, экспериментальное и теоретическое свидетельство, мира, который к тому же отождествлять с истинной реальностью --- то, что отказываются делать все феноменологи. Сама идея определения реально реальной реальности как неизвестной причины нашего опыта и полагание, скажем, объектов восприятия лишь знаками отдельных, скрытых объектов, чья истинная природа должна оставаться неизвестной и никогда не может быть постигнута --- не более чем мифологизация, считает Гуссерль. Вместо такого определения объективной реальности как существующей самой по себе, вместо разделения явленного и действительного и затем настаивания на действительно важности исследования лишь последнего --- вместо этого следует принять, что объективность включает как субъективность, так и интерсубъективность. Для Гуссерля (интер)субъективность --- не антипод объективности, и не препятствие для научного познания, он принимает концепцию объективности и реальности как данного, а не скрытого по ту сторону данного, тем самым избегая скептицизма касательно реального мира. Объективность для Гуссерля определяется не в терминах независимости от разума, а скорее в отношении общности разумов. Объективность --- результат интерсубъективного процесса, то, о чем мы можем договориться в результате исследования.

Настойчивое утверждение реальности лишь могущего быть известного естествознанию не только не позволяет должным образом рассмотреть и решить философский вопрос о реальном и объективном --- оно самоподрывно, поскольку не в состоянии адекватно объяснить экспериментальные и когнитивные достижения, делающие натурализм вообще возможным.

Наука часто подается как попытка объективно, т.е. от третьего лица, описать реальность. Стремление к объективности, конечно, похвально, но не следует забывать, что всякая объективность, всякое понимание и теоретическое моделирование, предполагают первоначально перспективу от первого лица, и в этом смысле иллюзорна возможность науки дать абсолютное описание реальности, свободное от всех теоретических и экспериментальных точек зрения. Наука --- особое отношение к миру, особая теоретическая модификация естественного отношения, и эта теоретическая установка не с потолка взята, а имеет свои предпосылки и генезис. Это традиция, культурное образование, это знание, разделяемое сообществом субъектов, переживающих опыт --- знание, предполагающее триангуляцию точек зрения или перспектив. Поэтому противопоставление перспективы от первого лица и от третьего вводит в заблуждение --- оно призывает нас забыть о создании отчетов от третьего лица сообществом сознательных субъектов, созерцающих и действующих от первого лица. ''Чистого'' вида от третьего лица не существует, так же как и вида из ниоткуда. Это, впрочем, не значит, что перспективы от третьего лица совсем нет, просто она все еще откуда-то, она все еще один из взглядов на мир. Наука уходит корнями в проживаемый мир, опирается на идеи донаучной сферы и осуществляется воплощенными субъектами, и стремясь понять ее эффективность и пределы, мы должны исследовать формы интенциональности, используемые познающими субъектами. Настройки счетчика, компьютерные распечатки, рентгеновские снимки и т.п. --- все это бессмысленно без сознательных субъектов, способных это интерпретировать и обсуждать. Сознание, таким образом --- условие объективности и поиска научных знаний гораздо более важное, чем, скажем, микроскопы и телескопы, и уж точно не помеха или препятствие. Феноменология стремится не предложить конкурирующее научное объяснение человеческих существ, а прояснить нашу научную практику, ее рациональность и достижения посредством детального анализа видов используемой познающими субъектами интенциональности. Каким образом теоретическая установка, которую мы используем, занимаясь наукой, возникает из нашего дотеоретического бытия-в-мире и изменяет его?

\chapter{От поверхностной к глубинной феноменологии}

Один из способов оценить развитие феноменологии --- рассматривать ее как постоянное расширение, углубление и усложнение основного феноменологического анализа, который можно найти в ранних работах Гуссерля. Сам Гуссерль в конечном итоге стал называть ту феноменологию, о которой изначально говорил, статической. Рассматривая некоторые из ранних формирующих анализов перцептивной интенциональности и отношений между восприятием и воображением, можно заметить, что все они изучают интенциональную корреляцию без учета генезиса и историчности, тип объекта и намеренного действия считаются легкодоступными. Но впоследствии Гуссерль пришел к пониманию субъекта интенциональности не как лишь формального принципа конституирования, но ''мертвым полюсом идентичности'', и то же справедливо для объектной стороны. Гуссерль продолжил исследование установления моделей понимания и их влияния на последующий опыт. В процессе ''оседания'' (sedimentation) опыт оставляет в нас свой след и тем самым способствует формированию когнитивных схем и разнообразных форм опасений и ожиданий, направляющих и мотивирующих последующие переживания. Определенные типы интенциональности --- например, долингвистический опыт --- обуславливают более поздние и сложные --- например, концептуальные суждения, и Гуссерль поставил задачу развития ''генетической феноменологии'' --- задачу исследования временного становления различных форм интенциональности и прослеживания форм объективности от более высоких порядков до более низких. Однако сфера генетической феноменологии по-прежнему ограничивалась эмпирической жизнью индивидуального эго. На последнем этапе Гуссерль отважился заняться так называемой генеративной феноменологией: для исследования конститутивной роли традиции и истории фокус был расширен. Каким образом достижения предыдущих поколений влияют на наш индивидуальный опыт? Как выразился Гуссерль в рукописи двадцатых годов: ''все мое собственное основано частично на традициях моих предков и частично --- на традициях моих современников''.

\section{Генеративность и традиция}

Один из способов проследить это развитие в собственных трудах Гуссерля --- повышенное внимание к переплетениям тем воплощенности, временности и социальности. Гуссерль, как уже указывалось, утверждал участие тела в восприятии пространственных объектов и во взаимодействии с ними, мир дан нам как исследуемый телесно, и тело раскрывается в своем исследовании мира. Уже на раннем этапе Гуссерль также осознал важность темпоральности и отстаивал неполноту исследования интенциональности при игнорировании временного измерения интенциональных действий и объектов. Он был одним из первых философов, всесторонее и систематически использовавших и обсуждавших понятие интерсубъективности и работал над эмпатией более трех десятилетий, в конечном итоге заявив, что феноменология должна развиться из ''эгологической'' в ''социологическую''.

В последних произведениях Гуссерля темы воплощенности, интерсубъективности и временности объединены и мыслятся вместе. Есть также диахроническое измерение интерсубъективности. Гуссерль в конечном счете считал, что рождение субъекта в живой традиции имеет конститутивные последствия, и дело не только в нахождении в пронизанном ссылками на других и уже другими наделенном смыслом, мире, или в моем понимании мира (и себя) через традиционные, унаследованные лингвистические условности --- сам смысл мира для меня таков, что он имеет свое начало вне меня, в историческом прошлом. Как пишет Гуссерль в ''Кризисе'', бытие, включенное в ''единый поток исторического развития'' --- в порождающую связь рождения и смерти --- так же неразрывно принадлежит Я, как и его временная форма.

Включение Гуссерлем в конечном итоге в свой трансцендентальный анализ таких тем как воплощенность, историчность и интерсубъективность --- возможно, одна из причин, почему Мерло-Понти в предисловии к ''Феноменологии восприятия'' написал: ''Трансцендентальное Гуссерля --- не то же, что трансцендентальное Канта''. Поздние идеи Гуссерля о взаимосвязи между ''я'', другими и миром с самого начала также развиваются Хайдеггером и Мерло-Понти, и они настаивают на принадлежности ''я'', мира и других друг другу, их взаимной явленности и возможности быть понятыми только в взаимосвязи. В одной из первых лекций Хайдеггер описывает проживаемый мир как взаимопроникновение трех областей: окружающего мира, с-мира (with-world) и я-мира (self-world), и утверждает, что Dasein как бытие-в-мире всегда уже бытие-с (Mitsein) другими. В более поздних лекциях 1927 года он пишет:

\smallskip
\leftskip1cm\relax
\rightskip1cm\relax

Dasein не просто поначалу является лишь обитанием среди вещей, чтобы затем время от времени обнаруживать среди этих вещей существ со своим собственным бытием
 --- напротив, как занятое собой бытие, Dasein столь же изначально бытие-с другими и бытие-среди внутримировых существ.

\leftskip0cm\relax
\rightskip0cm\relax
\smallskip

Мерло-Понти же утверждает, что субъективность по своей сути ориентирована и открыта тому, чем она является, и что именно в этой открытости она раскрывается самой себе. Феноменологическая рефлексия, значит, раскрывает не замкнутый в себе разум --- не чистое внутреннее самоприсутствие --- но открытость к инаковости, движение экстериоризации и постоянного самопревосхождения. Будучи представленными в миру, мы представлены себе, и будучи представлены себе, можем осознавать мир. Мерло-Понти, таким образом, настаивает на раскрытие преемственности между интерсубъективной жизнью и миром через феноменологическое описание, а не через раскрытие недоступных и самодостаточных субъективностей. Субъект представлен и себе, и миру, и другим не вопреки, но именно посредством телесности и историчности, и, согласно Мерло-Понти, такая концепция переплетения себя, других и мира задает концепцию трансцендентальной философии:

\smallskip
\leftskip1cm\relax
\rightskip1cm\relax

Как могут границы трансцендентального и эмпирического не быть нечеткими, когда вместе с другим человеком все, что он во мне видит --- вся моя фактичность --- реинтегрируется в субъективность или, по крайней мере, постулируется как обязательный элемент ее определения. Трансцендентальное так нисходит в историю, или, говоря иначе: историческое более не внешнее отношение между двумя или более абсолютно автономными субъектами, но имеет внутренний аспект и неотъемлимо от самого их определения. Они более не осознают себя субъектами лишь по отношению к себе, но также и по отношению друг к другу.

\leftskip0cm\relax
\rightskip0cm\relax
\smallskip

Чем феноменологическая работа Мерло-Понти и Хайдеггера отличается от таковой Гуссерля? Можно рассматривать оба подхода как реализацию уже имевшихся у Гуссерля идей в более радикальном ключе.

Так, Мерло-Понти придает большее, чем Гуссерль, значение роли воплощения и фактичности, то есть контингентной ситуированности нашего существования, а также идет дальше в своей попытке переосмыслить традиционное разделение между трансцендентальным и эмпирическим и между разумом и миром. Это очевидно уже в первой его крупной работе ''Структура поведения'' (1942), где он активно обращается к эмпирической науке и на последней странице призывает к переопределению трансцендентальной философии: вместо выбора между научным объяснением и феноменологической рефлексией, Мерло-Понти призывает обращать внимание на живые отношения между сознанием и природой и искать измерение за пределами как объективизма, так и субъективизма.

Хайдеггер же внимательнее читает историю философии, нежели Гуссерль, а также больше подчеркивает влияние традиции на современное мышление. Для Хайдеггера одной из важных задач феноменологии является раскрытие и деконструкция некоторых метафизических концепций, на протяжении веков формировавших и ограничивавших философское мышление.

\end{document}
